\chapter{Reading a Chest X-Ray}

\label{ch:read_cxr}

Interpretation of a chest radiograph (CXR) is a fundamental clinical skill required across all medical specialties. A haphazard look at a CXR often leads to missed diagnoses. To ensure accuracy, clinicians must apply a systematic approach every time.

---

\section*{1. Technical Quality Assessment: The RIPE Criteria}
Before attempting to interpret any pathological findings, you must evaluate the technical quality of the image. An inadequate projection can easily mimic or obscure serious diseases. Use the mnemonic \textbf{RIPE}:

\begin{itemize}
    \item \textbf{Rotation}: Inspect the clavicular heads. The medial ends of the clavicles should be equidistant from the spinous process of the thoracic vertebra at that level. If the patient is rotated, the mediastinum and heart borders will appear distorted, potentially mimicking cardiomegaly or mediastinal widening.
    \item \textbf{Inspiration}: A good inspiratory effort is crucial. Count the visible ribs above the diaphragm. You should be able to see \textbf{9 to 10 posterior ribs} or \textbf{6 anterior ribs} projecting above the right hemidiaphragm. Poor inspiration compresses the lung bases, mimicking basal consolidation, and falsely enlarges the cardiac silhouette.
    \item \textbf{Projection}: Determine if the film is Posteroanterior (PA) or Anteroposterior (AP).
    \begin{itemize}
        \item \textbf{PA view}: The standard view. The patient stands with the chest against the detector. Scapulae are rotated laterally out of the lung fields. The heart is closer to the detector, minimizing magnification.
        \item \textbf{AP view}: Typically performed portably for critically ill patients. Scapulae overlap the lung fields, and the heart is magnified because it is farther from the detector. Do not diagnose cardiomegaly on an AP film alone.
    \end{itemize}
    \item \textbf{Exposure (Penetration)}: Assess the visibility of the thoracic spine. You should just be able to visualize the vertebral body outlines through the lower part of the cardiac silhouette. 
    \begin{itemize}
        \item \textbf{Under-penetrated (too white)}: Left hemidiaphragm and lung bases are obscured.
        \item \textbf{Over-penetrated (too dark)}: Lung markings are burnt out, potentially masking subtle nodules or pneumothoraces.
    \end{itemize}
\end{itemize}

---

\section*{2. Systematic Interpretation: The ABCDE Method}
Once technical adequacy is confirmed, interpret the image systematically from the center outward:

\subsection*{A --- Airway}
\begin{itemize}
    \item \textbf{Trachea}: Ensure the trachea is midline. Note any deviation. Tracheal deviation is a crucial sign:
    \begin{itemize}
        \item \textit{Pulled toward the pathology}: Ipsilateral lung collapse or pneumonectomy.
        \item \textit{Pushed away from the pathology}: Tension pneumothorax or massive pleural effusion.
    \end{itemize}
    \item \textbf{Carina and Bronchi}: Identify the bifurcation (carina) and trace the right and left mainstem bronchi. The angle of the carina should normally be between 60 and 100 degrees; widening can indicate subcarinal lymphadenopathy or left atrial enlargement.
\end{itemize}

\subsection*{B --- Breathing (Lungs and Pleura)}
Divide each lung into three zones (upper, middle, lower) and compare them side-to-side.
\begin{itemize}
    \item \textbf{Lung Parenchyma}: Check for symmetry in radiolucency. Look for:
    \begin{itemize}
        \item \textbf{Consolidation}: Alveoli filled with fluid, pus, or blood, appearing as homogeneous white opacities.
        \item \textbf{Interstitial markings}: Reticular or nodular patterns indicating pulmonary edema or fibrosis.
        \item \textbf{Nodules/Masses}: Well-circumscribed lesions requiring evaluation for malignancy.
    \end{itemize}
    \item \textbf{Pleural Borders}: Trace the edges of the lungs. A thin, sharp black line at the periphery with a complete absence of distal bronchovascular markings indicates a \textbf{pneumothorax}.
    \item \textbf{Pleural Spaces}: Check for \textbf{pleural effusion}. A blunting of the costophrenic angle indicates fluid accumulation (requires ~200 mL on a PA view to be visible).
\end{itemize}

\subsection*{C --- Circulation (Heart and Mediastinum)}
\begin{itemize}
    \item \textbf{Heart Size}: Evaluate the cardiothoracic ratio (CTR) on a PA projection. The CTR should be \textbf{< 50\%} in healthy adults. A CTR > 50\% indicates cardiomegaly.
    \item \textbf{Heart Borders}: Identify the right atrium border (right side) and left ventricle border (left side).
    \item \textbf{Mediastinal Contour}: Inspect the aortic knob, pulmonary trunk, and left atrial appendage. Widening of the superior mediastinum (> 8 cm on PA view) warrants evaluation for aortic dissection, mediastinal masses, or lymphadenopathy.
    \item \textbf{Hila}: The left hilum is normally slightly higher than the right hilum (by 1--2 cm). Compare hilar density and shape; bulky, lobulated hila indicate lymphadenopathy or pulmonary hypertension.
\end{itemize}

\subsection*{D --- Diaphragm}
\begin{itemize}
    \item \textbf{Diaphragmatic Domes}: The right hemidiaphragm is normally slightly higher than the left (by 1.5--2 cm) due to the liver.
    \item \textbf{Flatness}: Flattened diaphragms indicate pulmonary hyperinflation (chronic obstructive pulmonary disease).
    \item \textbf{Free Air}: Inspect below the right hemidiaphragm. A thin crescent of air (crescent sign) represents \textbf{pneumoperitoneum}, indicating a perforated abdominal viscus (a surgical emergency).
\end{itemize}

\subsection*{E --- Everything Else (Bones, Soft Tissues, and Lines)}
\begin{itemize}
    \item \textbf{Bones}: Trace the clavicles, scapulae, humeri, and ribs. Check for rib fractures, osteolytic/osteoblastic lesions, and spinal alignment.
    \item \textbf{Soft Tissues}: Look for subcutaneous air (subcutaneous emphysema) which feels like crackling paper and presents as dark streaks in the neck or chest wall.
    \item \textbf{Medical Devices}: Verify the positioning of all support lines and tubes:
    \begin{itemize}
        \item \textbf{Endotracheal (ET) Tube}: The tip should lie 3--5 cm above the carina (approximately at the level of the clavicular heads).
        \item \textbf{Nasogastric (NG) Tube}: The tube must descend midline, pass below the diaphragm, and the tip must lie in the stomach stomach gastric bubble.
        \item \textbf{Central Venous Pressure (CVP) Line}: The tip should lie in the superior vena cava, just above its junction with the right atrium.
    \end{itemize}
\end{itemize}

---

\section*{3. Hallmark Radiological Signs}
Understanding these classic imaging signs is key to forming a differential diagnosis:

\begin{itemize}
    \item \textbf{The Silhouette Sign}: The loss of a normal anatomical interface (such as the heart border or diaphragm) caused by an adjacent fluid collection or consolidation of similar density.
    \begin{itemize}
        \item Loss of the right heart border = Right middle lobe consolidation.
        \item Loss of the left heart border = Left lingular consolidation.
        \item Loss of the hemidiaphragm border = Lower lobe consolidation.
    \end{itemize}
    \item \textbf{Air Bronchograms}: The visualization of air-filled, dark bronchi standing out against surrounding opacified (white) alveoli. This indicates airspace disease (e.g., lobar pneumonia) rather than atelectasis.
    \item \textbf{Deep Sulcus Sign}: A very deep, lucent costophrenic angle seen on a supine chest radiograph, indicating a pneumothorax in patients lying down.
\end{itemize}

---

\section*{4. Structured Reporting Examples}

To maintain professional standards, radiologists use structured reporting. Here are template examples:

\subsection*{Normal Chest X-Ray Report}
\begin{testbox}[Indications:]
Chest pain and dyspnea. Rule out pneumonia.
\end{testbox}
\begin{refbox}[Findings:]
\textbf{Airway}: Trachea is midline. Normal carinal angle.\\
\textbf{Breathing}: Lungs are clear. No focal consolidation, pleural effusion, or pneumothorax. Lungs are normally inflated.\\
\textbf{Circulation}: Cardiomediastinal silhouette is within normal limits. Cardiothoracic ratio is less than 50\%. Hila are normal in position and configuration.\\
\textbf{Diaphragm}: Hemidiaphragms are smooth and sharply defined. Costophrenic angles are clear. No free air below the diaphragm.\\
\textbf{Bones \& Soft Tissues}: Osseous structures and chest wall soft tissues are unremarkable. No foreign bodies or indwelling lines are present.
\end{refbox}
\begin{pearlbox}[Impression:]
No acute cardiopulmonary disease.
\end{pearlbox}

\subsection*{Abnormal Chest X-Ray Report (Lobar Pneumonia)}
\begin{testbox}[Indications:]
Fever, productive cough, and tachypnea.
\end{testbox}
\begin{refbox}[Findings:]
\textbf{Airway}: Trachea is midline. Carinal bifurcation is clear.\\
\textbf{Breathing}: Homogeneous consolidation with prominent air bronchograms is noted in the right middle lobe. This consolidation obscures the right heart border (silhouette sign). Left lung is clear. No pleural effusion or pneumothorax is seen.\\
\textbf{Circulation}: Heart size is normal. Superior mediastinum is not widened.\\
\textbf{Diaphragm}: Hemidiaphragms are intact. Costophrenic angles are sharp.\\
\textbf{Bones \& Soft Tissues}: Intact thoracic skeleton. Soft tissues are normal.
\end{refbox}
\begin{pearlbox}[Impression:]
Right middle lobe consolidation, highly suggestive of acute lobar pneumonia. Recommend clinical correlation and follow-up radiograph post-antibiotic treatment to ensure clearance.
\end{pearlbox}

---

\section*{5. Common Pitfalls and How to Avoid Them}
\begin{itemize}
    \item \textbf{Misinterpreting Skin Folds as a Pneumothorax}: In thin patients or elderly patients, loose skin folds can fold against the detector, mimicking a pneumothorax line. To differentiate: skin folds usually extend beyond the thoracic cage, and you will see lung markings present distal to the fold line (which is not the case in a true pneumothorax).
    \item \textbf{Mistaking Expiratory Films for Pathology}: An expiratory film will cause the heart to appear enlarged and the lung bases to appear opacified. Always count the ribs to verify inspiration before diagnosing cardiomegaly or basal atelectasis/fluid.
\end{itemize}

---

\section*{6. AI-Assisted Interpretation}
Modern clinical systems utilize Artificial Intelligence (AI) algorithms as concurrent readers. AI is highly sensitive for:
\begin{itemize}
    \item Rapid triage of life-threatening conditions (detecting a pneumothorax or tube misplacement in under 2 minutes).
    \item Automated detection and tracking of pulmonary nodules.
    \item Calculating cardiothoracic ratios automatically.
\end{itemize}
However, AI can throw false positives on rotated films or vascular congestion. Thus, all AI findings must be verified by clinical review.

\clearpage
